\documentclass[12pt]{article}
\usepackage[utf8]{inputenc}
\usepackage[T1]{fontenc}
\usepackage{hyperref}
\title{Puwa Yi Language Resources for K-12 Curriculum Development}
\author{Compiled by Claude Sonnet 5.5 (claude-sonnet-5-5) for Andr\'as Kornai}
\date{}
\begin{document}
\maketitle

\section{Introduction}
The code ypw (Puwa Yi) is retired. According to the Glottolog bookkeeping entry, it was split into three languages: Hlepho Phowa [yhl], Labo Phowa [ypb] and Ani Phowa [ypn]. Wikipedia describes Phowa as a dialect cluster of Loloish languages spoken by the Phula people in southern Yunnan, China. Reported figures are roughly 36,000 (Hlepho), 17,000 (Labo) and 10,000 (Ani), dated 2011. Sources disagree on how the varieties should be classified. The searches found no curriculum materials specific to ypw.

\section{Primary Educational Resources}
No primers, textbooks or school programmes for Puwa Yi turned up in the searches.

\section{Linguistic and Reference Materials}
\begin{itemize}
\item Glottolog bookkeeping entry for Puwa Yi, which gives the successor codes: \url{https://glottolog.org/resource/languoid/id/puwa1234}
\item Wikipedia page on the Phowa languages (retrieved via the ISO 639 yhl entry): \url{https://wikipedia.com/wiki/ISO_639:yhl}
\end{itemize}

\section{Religious and Community Resources}
\begin{itemize}
\item Asia Harvest people-group profile of the Puwa in Yunnan. It is an older missionary source and its figures may be dated: \url{https://people-groups.asiaharvest.org/people-group-profiles/yunnan/puwa.pdf}
\end{itemize}

\section{Digital and Audio Resources}
\begin{itemize}
\item Global Recordings Network entry titled ``Yi: Puwa in Kaiyuan''. I did not verify which variety it covers: \url{https://globalrecordings.net/en/language/6073}
\end{itemize}

\section{Curriculum Development Context}
Curriculum work should target the successor languages (Hlepho, Labo and Ani Phowa), not ypw. The resources above are identification and background material only. No native-language school materials were found.
\end{document}